\documentclass[11pt]{article}
\usepackage[margin=1in]{geometry}
\usepackage{natbib}
\usepackage{booktabs}
\usepackage{amsmath}
\usepackage{graphicx}
\usepackage{hyperref}
\usepackage{authblk}

\title{Cryo-EM Structure and B-Factor Refinement with Ensemble Representation}

\author[1]{Authors}
\affil[1]{Bijvoet Center for Biomolecular Research, Faculty of Science -- Chemistry, Utrecht University, Utrecht, The Netherlands}

\date{}

\begin{document}
\maketitle

\begin{abstract}
Cryo-EM experiments produce images of macromolecular assemblies that are combined to produce three-dimensional density maps. It is common to fit atomic models of the contained molecules to interpret those maps, followed by a density-guided refinement. Here, we propose TEMPy-REFF, a novel method for atomic structure refinement in cryo-EM density maps. By representing the atomic positions as components of a mixture model, their variances as B-factors, and a model ensemble description, we significantly improve the fit to the map compared to what is currently achievable with state-of-the-art methods. We validate our method on a large benchmark of 366 cryo-EM maps from EMDB at 1.8--7.1~\AA{} resolution and their corresponding PDB assembly models. We also show that our approach can provide newly modelled regions in EMDB-deposited maps by combining it with AlphaFold-Multimer. Finally, our method provides a natural interpretation of maps into components, allowing us to accurately create composite maps.
\end{abstract}

\section{Introduction}

Cryo-electron microscopy (cryo-EM) can resolve the structure of biomolecules at an ever-improving resolution \citep{nakane2018}. Larger complexes can now be visualised as three-dimensional density maps at near-atomic resolutions, and in various conformations. The interpretation of those maps often hinges on fitting atomic models of the different macromolecules present in the complex \citep{joseph2017,miyashita2012,topf2008}. This procedure is often difficult and requires the user to provide accurate models and a well-estimated resolution, which can vary at different parts of the map. Pre-existing experimental or predicted atomic models may be in a different conformation and converging to a well-fitted one may require significant sampling.

Several methods are in common use for this task: molecular dynamics (MD)-based approaches \citep{igaev2019,hoffmann2015}, real-space refinement in Phenix \citep{afonine2018,liebschner2019}, flexible fitting through normal mode analysis \citep{lopez2013,tama2004}, Rosetta-based refinement \citep{dimaio2016}, and interactive model building tools such as Coot \citep{emsley2010} and ISOLDE \citep{croll2018}. Gaussian mixture model (GMM) representations have been used for map fitting and representation \citep{kawabata2008,kawabata2018}. However, B-factor refinement --- which encodes the local flexibility and quality of fit --- has received comparatively little attention in the cryo-EM context.

Here, we present TEMPy-REFF, a method that frames atomic structure refinement as an Expectation-Maximisation (EM) algorithm operating on a Gaussian mixture model. By treating atomic positions as mixture components and their B-factors as variances, we achieve significant improvements in map-model correlation. We further generate structural ensembles based on refined B-factors and demonstrate applications in map segmentation and composition.

\section{Methods}

\subsection{Algorithm overview}

Given a reference experimental 3D map $M_E$ with an estimated resolution $R$, we attempt to improve the fit of a model comprised of $N$ atoms with estimated coordinates $\mathbf{X}$, B-factors $\mathbf{B}$, and occupancies (set to 1). The algorithm iterates over the following steps:
\begin{enumerate}
    \item Compute conformational energy and forces for a given force field.
    \item Compute a simulated density map $M_S$.
    \item Compute map-derived forces.
    \item Update coordinates $\mathbf{X}$.
    \item Update B-factors based on the experimental cryo-EM map $M_E$.
\end{enumerate}

\subsection{Simulated cryo-EM map}

From the positions $\mathbf{X}$ and B-factors $\mathbf{B}$, we compute a simulated map $M_S$ by modelling each atom as a Gaussian:
\begin{equation}
    M_S(\mathbf{x}) = \sum_{i=1}^{N} N_i \cdot \mathcal{N}(\mathbf{x}; \mathbf{x}_i, B_i) + k_{bg}
\end{equation}
where $N_i$ is the atomic number, $\mathbf{x}_i$ the Cartesian coordinates, $B_i$ the B-factor, and $k_{bg}$ a uniform background noise parameter.

\subsection{Ensemble generation}

To represent the variety of conformations compatible with the map, we generate a structural ensemble based on refined B-factors. Models are generated by randomly perturbing atomic positions according to their B-factors, followed by local l-BFGS minimisation \citep{liu1989} using OpenMM \citep{eastman2017}.

\subsection{Map segmentation and composition}

The mixture model representation naturally assigns a ``responsibility'' to each atom for each voxel of the map. This allows segmentation of the map into submaps corresponding to distinct chains or subunits, as well as composition of focused maps into a single composite map with optimal weighting.

\subsection{Benchmark dataset}

We validated our method on the CERES+ benchmark of 388 protein complexes \citep{ceres2021}, based on the recently released CERES database with additional cases, containing 366 cryo-EM maps at resolutions $\leq$5~\AA{} from EMDB and their corresponding PDB assembly models.

\section{Results}

\subsection{B-factor refinement converges rapidly and is robust}

The B-factor converges quite rapidly, usually in 10--11 iterations or less, with minimal changes observed afterwards (Extended Data Fig.~S1, S2). This was demonstrated on examples including Faba bean necrotic stunt virus (FBNSV, EMD-10097) and the SARS-CoV-2 RNA-dependent RNA polymerase (EMD-30127). The B-factor assignment is highly robust: recovered B-factors are near-identical after a few iterations, with deviations of $\pm$0.87~\AA$^2$ over 366 cases regardless of the initial B-factor values (Extended Data Fig.~S3, S4, S5).

\subsection{Ensemble generation improves correlation}

To more accurately represent conformational diversity, we generate structural ensembles based on refined B-factors (Figure~2). For rotavirus VP6 (EMD-6272) at 2.6~\AA, the ensemble provides improved correlation compared to a single model (Figure~2a--c). Different residues show distinct ensemble spread: arginine 107 shows a more widespread ensemble with diffuse side-chain density, while arginine 217 shows a highly constrained ensemble with well-defined density (Figure~2d).

\subsection{Benchmarking structure refinement}

\begin{table}[htbp]
\centering
\caption{Change in CCC upon B-factor refinement in TEMPy-REFF. Initial values are listed for a uniform B-factor of 10; final values are listed after 5 B-factor refinement steps. Atomic coordinates are not changed during this process.}
\label{tab:bfactor}
\begin{tabular}{lcc}
\toprule
Resolution band & Initial CCC & Final CCC \\
\midrule
1.8--3.0 \AA & 0.52 $\pm$ 0.12 & 0.74 $\pm$ 0.09 \\
3.0--4.0 \AA & 0.48 $\pm$ 0.11 & 0.69 $\pm$ 0.10 \\
4.0--5.0 \AA & 0.43 $\pm$ 0.13 & 0.62 $\pm$ 0.12 \\
\bottomrule
\end{tabular}
\end{table}

We benchmarked TEMPy-REFF against the CERES dataset \citep{ceres2021} (Table~\ref{tab:bfactor}, Figure~3). The improvements in CCC scores are statistically significant compared with results obtained for the deposited and CERES models using paired-sample T-tests: average improvements of 0.197 with a $p$-value of 0.0006 compared to EMDB models and 0.150 with a $p$-value of 0.003 compared to CERES models. These improvements hold across all resolution bands, though they become more pronounced at lower resolutions where the deposited and CERES models show lower CCC values (Figure~3a).

\begin{table}[htbp]
\centering
\caption{Change in scores upon refinement, comparison against CERES.}
\label{tab:scores}
\begin{tabular}{lccc}
\toprule
Metric & Deposited & CERES & TEMPy-REFF \\
\midrule
CCC (mean) & 0.64 & 0.69 & 0.84 \\
SMOCf (mean) & 0.71 & 0.74 & 0.82 \\
MolProbity & 2.8 & 2.5 & 2.3 \\
CaBLAM outliers (\%) & 4.2 & 3.8 & 3.1 \\
\bottomrule
\end{tabular}
\end{table}

Overall local measures of fit quality (SMOC) show improvements with TEMPy-REFF models. In the example of glutamate dehydrogenase (PDB ID 5K12, EMD-8194), improvements are apparent in most regions of the modelled structure (Figure~3b). CaBLAM scores also show improvement (Table~\ref{tab:scores}).

\begin{table}[htbp]
\centering
\caption{Speed benchmark: time for a refinement run for a given system.}
\label{tab:speed}
\begin{tabular}{lcc}
\toprule
System & Atoms & Time (min) \\
\midrule
Rotavirus VP6 & 3,200 & 12 \\
Glutamate dehydrogenase & 13,000 & 45 \\
RNA Pol III & 35,000 & 120 \\
\bottomrule
\end{tabular}
\end{table}

\begin{table}[htbp]
\centering
\caption{Force field comparison: change in CCC after refinement using either AMBER or CHARMM.}
\label{tab:forcefield}
\begin{tabular}{lcc}
\toprule
Force field & Mean $\Delta$CCC & Ramachandran outliers (\%) \\
\midrule
AMBER (ff14SB) & +0.19 & 0.8 \\
CHARMM36 & +0.18 & 0.9 \\
\bottomrule
\end{tabular}
\end{table}

No significant bias is introduced by the choice of force field, as confirmed by testing both AMBER \citep{maier2015} and CHARMM36 \citep{best2012} (Table~\ref{tab:forcefield}, Extended Data Fig.~S9).

\subsection{Map segmentation and composition}

We demonstrate map segmentation on the 2.8~\AA{} resolution glycoprotein B-neutralizing antibody complex from Herpes Simplex Virus 1 (EMD-2124736), segmenting the map into 9 submaps corresponding to single chains (Figure~4a). For map composition, we combine 3 focused maps of the RNA polymerase II super elongation complex (EMD-12966 to EMD-12969) into a composite map (Figure~4b). Because the responsibility decays smoothly, there are no ``seams'' between segmented or composite maps.

\subsection{Case study I: Yeast RNA polymerase III elongation complex}

We refined the deposited model (PDB ID 5FJ8) in the 3.9~\AA{} map (EMD-3178). The deposited model had a correlation score of 0.82 and MolProbity score of 2.8. TEMPy-REFF refinement yielded a final CCC of 0.94 for the ensemble (0.83 for the single model). The average LoQFit score improved from 10.1~\AA{} to 9.0~\AA{} (Figure~5), with particularly notable improvements in lower-resolution peripheral chains where average LoQFit scores improved from 7.4 to 6.0~\AA.

\subsection{Case study II: Nucleosome-CHD4 complex}

We applied TEMPy-REFF to refine the model (PDB ID 6RYR) in map EMD-10058 (Figure~6a). The deposited map suffers from variable resolution (3--10~\AA), which affected the quality-of-fit. TEMPy-REFF refinement improved the fit, with regions of higher B-factor (lower local resolution) showing correspondingly higher ensemble variability (Figure~6b--c).

\subsection{Case study III: SARS-CoV-2 RNA polymerase with AlphaFold2}

To assess refinement of predicted models, we used the AlphaFold2 \citep{jumper2021} prediction for the SARS-CoV-2 polymerase (UNIPROT ID: P0DTD1, residues 4393--5324) and refined it into the 2.9~\AA{} map (EMD-30127). The refined model is highly similar to the deposited model (PDB ID 6M71) at most residue positions. Using SMOCf plots, some residues not present in the deposited structure can be placed into the map with fitting scores much greater than chance (Figure~6e--h).

\section{Discussion}

TEMPy-REFF represents a significant advance in cryo-EM model refinement by integrating B-factor refinement, ensemble generation, and mixture model-based map representation. The method achieves statistically significant improvements over both deposited models and state-of-the-art CERES re-refined models across a benchmark of 366 structures.

The B-factor refinement is remarkably robust, converging to near-identical values regardless of initialisation. This property, combined with the ensemble representation, provides a natural measure of local model uncertainty that correlates with local map resolution \citep{cardone2013}. The mixture model framework additionally enables seamless map segmentation and composition without artificial boundaries.

The ability to refine AlphaFold2 predictions into experimental maps, potentially recovering residues absent from deposited models, suggests TEMPy-REFF can serve as a bridge between structure prediction and experimental validation \citep{jumper2021,tunyasuvunakool2021,evans2022}. Future development will focus on integrating multi-body refinement capabilities and extending the method to handle larger and more heterogeneous assemblies.

\section{Data Availability}

TEMPy-REFF is available as part of the TEMPy package \citep{farabella2015}. All maps and models used in this study are publicly available from EMDB and PDB.

\bibliographystyle{plainnat}
\bibliography{references}

\end{document}
